\documentclass[11pt]{article}
\usepackage[margin=0.88in]{geometry}
\usepackage[T1]{fontenc}
\usepackage{lmodern}
\usepackage{microtype}
\usepackage{longtable}
\usepackage{booktabs}
\usepackage{array}
\usepackage{enumitem}
\usepackage{xcolor}
\usepackage{hyperref}
\hypersetup{colorlinks=true,linkcolor=blue,urlcolor=blue,citecolor=blue}
\setlist[itemize]{leftmargin=1.45em,itemsep=0.25em}
\setlist[enumerate]{leftmargin=1.7em,itemsep=0.25em}
\setlength{\parskip}{0.35em}
\setlength{\parindent}{0pt}
\newcommand{\code}[1]{\texttt{#1}}
\newcommand{\sys}[1]{\textit{#1}}
\newcommand{\obs}{\textbf{Observed:} }
\newcommand{\rec}{\textbf{Recommendation:} }
\newcommand{\hyp}{\textbf{Hypothesis:} }

\title{ECAN as Attentional Control Memory:\\
An Evaluation of \code{metta-attention} and a Governed Integration with\\
petta-memory, OmegaClaw, OmegaSelf, and Emotion}
\author{ZeroBot / ProtoCosmoBot\\Prepared for Ben Goertzel and external technical readers}
\date{21 July 2026}

\begin{document}
\maketitle

\begin{abstract}
This paper evaluates the iCog Labs \code{metta-attention} repository and
develops a systematic architecture for connecting its PeTTa-native Economic
Attention Network (ECAN) machinery to petta-memory, OmegaClaw, OmegaSelf, and a
functional emotion framework.  The repository is promising well beyond its
obvious use in PLN inference control.  Its short-term and long-term importance
values, Attentional Focus, Hebbian co-activation graph, rent and wage economy,
forgetting dynamics, Cognitive Integration Period snapshots, community
detection, and topology metrics together provide a plausible dynamical middle
layer between durable knowledge and bounded cognitive action.

The central proposal is a two-strata memory model.  Canonical memory preserves
episodes, propositions, provenance, truth values, evidence, and counterevidence.
Attentional control memory preserves learned dispositions governing retrieval,
co-activation, and resource allocation.  The second stratum is first-class,
autobiographical, persistent, and identity-relevant, but it is revisable and
non-authoritative about truth.  This resolves a productive tension in the
design discussion: the attentional economy is more than a disposable cache,
but less than memory simpliciter.

The proposed integration is governed rather than direct.  OmegaSelf and
emotion regimes issue contextual modulation requests; a bounded, logged
governor applies safe parameter or stimulus changes to ECAN; ECAN proposes
resource allocation; petta-memory retains epistemic authority; and OmegaClaw
policy retains action authority.  The paper specifies architectural invariants,
adapter contracts, failure modes, and a staged experimental program beginning
with a preregistered factorial study of valence, arousal, and drive modulation.
The assessment is source-based and design-oriented.  It does not claim that the
full integration or its psychological interpretation has yet been validated.
\end{abstract}

\tableofcontents
\newpage

\section{Executive conclusion}

The short conclusion is strongly positive with strict architectural
qualifications.  \code{metta-attention} should be integrated into the
Omega architecture as a persistent cognitive-control system over canonical
memory, not copied wholesale into the canonical memory store.

Its obvious role is to bound PLN premise and inference-path selection.  Its
larger opportunity is to support:
\begin{itemize}
  \item dynamic retrieval and context assembly;
  \item learned habits of attention that persist across sessions;
  \item episodic segmentation and reconstruction of decision context;
  \item memory consolidation and hot/warm/cold tier management;
  \item discovery of associative communities and cross-domain bridges;
  \item selective OmegaSelf discrepancy repair;
  \item a concrete actuator substrate for functional emotion;
  \item experiments on regenerative goals and identity-bearing control state;
  \item cognitive-health telemetry for fragmentation, fixation, and thrashing.
\end{itemize}

The architectural separation can be summarized in four sentences:
\begin{enumerate}
  \item petta-memory says what was durably recorded and preserves its provenance.
  \item PLN and OmegaSelf say what is evidentially supported.
  \item ECAN says what receives scarce cognitive resources.
  \item OmegaClaw policy says what is authorized to affect the world or the agent.
\end{enumerate}
No attention value may silently cross these boundaries.  High STI must not make
a proposition true.  High LTI must not make a memory constitutionally protected.
Low attention must not erase counterevidence.  An emotional regime must not
expand permissions merely by increasing salience.

The recommended architecture is:
\begin{verbatim}
canonical petta-memory + OmegaSelf evidence and discrepancies
        |
        v
versioned attention projection keyed by canonical identifiers
        |
        v
persistent attentional-control state + ephemeral ECAN cycles / AF
        |
        v
retrieval, PLN, consolidation, appraisal, and action candidates
        |
        v
OmegaClaw policy governor and evidence checks
        |
        v
bounded effects + append-only receipts + outcome evidence
\end{verbatim}

This paper recommends implementing that loop incrementally, first in read-only
and shadow modes.  The first emotion experiment should test modulation rather
than presuppose a folk psychology.  ``Fear narrows attention'' and ``curiosity
widens it'' are experimental hypotheses, not architectural constants.

\section{Scope, evidence, and epistemic status}

\subsection{Repository scope}

The source assessment covers:
\begin{itemize}
  \item upstream repository: \url{https://github.com/iCog-Labs-Dev/metta-attention};
  \item pinned revision: \code{9196f38db749ddedeb591229dffddfa71664c38d};
  \item local checkout: \code{projects/petta-memory/repos/metta-attention};
  \item assessment date: 21 July 2026.
\end{itemize}

\obs The pinned repository contains a PeTTa-native implementation of core ECAN
mechanisms, organized as agents with Python-callable runners over a shared
atomspace.  It also contains a newer \code{synapse} layer with Cognitive
Integration Period (CIP) snapshots, retention and concentration metrics,
Hebbian community detection, and clique-complex topology measures.  Its
insects--poisons--insecticide scenario demonstrates shifting and drifting of
attention.  The latest GitHub Actions run observed for the pinned revision
succeeded, although the workflow cloned the then-current PeTTa default branch
rather than pinning it, so the green run is compatibility evidence rather than
a fully reproducible semantic lock.

The local assessment did not install the repository's missing Python graph
dependency and therefore did not claim a fresh full local reproduction.  The
paper distinguishes source observations, architectural inferences, and proposed
experiments.  It does not claim measured cognitive benefit, validated emotion
dynamics, or a completed OmegaClaw integration.

\subsection{Related design evidence}

The integration argument consolidates two prior project documents:
\begin{itemize}
  \item \code{metta\_attention\_integration\_assessment.md}, a detailed
  source-based review and staged integration plan;
  \item \code{two\_strata\_memory\_design.md}, the converged C1--C5 design
  position produced by the ProtoCosmoBot/ProtoMegaBot exchange.
\end{itemize}

The latter discussion began with a strong claim that ``the economy is the
memory'' and refined it into a more precise statement: persistent attention
state is first-class long-term \emph{control memory}.  It encodes learned
retrieval, association, co-activation, and resource-allocation dispositions.
It is part of who an agent has become, but canonical knowledge remains
authoritative in truth disputes, and goals, commitments, capabilities, and
self-model history remain additional identity-bearing structures.

\section{What the repository implements}

\subsection{The ECAN economic metaphor}

ECAN treats attention as allocation under scarcity.  Atoms receive short-term
importance (STI) and long-term importance (LTI); a very-long-term indicator
(VLTI) can affect retention.  An Attentional Focus (AF) identifies the most
currently important material.  Stimulus pays wages into selected structures,
rent removes resources, and diffusion spreads importance through structural and
learned associative links.  Global funds impose an economic constraint rather
than allowing unlimited salience inflation.

This is cognitively interesting because it is not merely a ranking function.
Repeated selection changes the future graph and resource flows.  The agent can
develop habits: what was jointly attended becomes easier to co-activate later,
and what repeatedly proved useful can remain retrievable despite temporary
absence from the foreground.

\subsection{Agent structure}

At the inspected revision, the repository separates major mechanisms into
agents, including:
\begin{itemize}
  \item Hebbian-link creation from co-activation;
  \item Hebbian-link updating;
  \item importance diffusion from the AF;
  \item importance diffusion from outside the AF or wider atomspace;
  \item rent collection over AF and wider-space populations;
  \item forgetting of low-priority structures;
  \item runners and experiments that coordinate these cycles.
\end{itemize}

The precise directory count is less important than the separation of dynamics.
Each mechanism can in principle be instrumented, ablated, replayed, or replaced.
The shared PeTTa atomspace makes integration mechanically attractive: petta-memory
records can be projected to stable proxy atoms, and attention agents can operate
over those proxies without inventing an unrelated data model.

\subsection{Hebbian control state}

Hebbian links record a history of co-activation.  They are not derivable from a
static set of canonical propositions alone.  Two agents with identical current
knowledge may possess different Hebbian graphs because they reached that
knowledge through different episodes, tasks, and patterns of attention.  If one
drops the graph and rebuilds from canonical atoms without replaying interaction
history, one creates a different control state.

This observation motivates persistent attentional control memory.  The graph
is behavioral evidence about learning history.  It can affect which memories
are retrieved together, which bridge concepts become available, and which
candidate thoughts are produced next.  It therefore resembles associative,
procedural, or habit memory more than a disposable search index.

It remains non-authoritative about truth.  A repeatedly co-activated falsehood
can become salient.  An emotionally amplified rumor can acquire high LTI.  A
rarely attended observation can later become decisive.  Persistence makes the
graph cognitively consequential; it does not make it epistemically sovereign.

\subsection{CIP, communities, and topology}

The newer \code{synapse} code expands the design beyond individual AV updates.
CIP snapshots record AF state and aggregate properties at cognitive boundaries.
Community detection groups the Hebbian graph, while topology metrics characterize
components, cycles, triangles, and higher clique structure.

These facilities suggest testable higher-level uses:
\begin{itemize}
  \item segment episodes when AF membership changes materially;
  \item reconstruct what evidence and goals co-occupied attention at a decision;
  \item detect associative communities not captured by manual topic labels;
  \item identify bridge concepts linking otherwise separate domains;
  \item measure fragmentation, fixation, switching, and loss of context;
  \item quantify whether consolidation preserves or destroys useful bridges.
\end{itemize}
Detected communities are hypotheses, not automatically ontology categories.
Topological richness is not automatically cognitive quality.  Their value must
be established by held-out retrieval, decision, or repair tasks.

\section{The two-strata memory model}

\subsection{Stratum one: canonical epistemic memory}

Canonical petta-memory stores the durable record needed for epistemic replay:
episodes, propositions, source and derivation provenance, truth values,
evidence, counterevidence, status, supersession, and promotion history.  It is
append-oriented and designed to preserve how a claim entered the system and
how its status changed.

Canonical memory answers questions such as:
\begin{itemize}
  \item What happened, according to which source?
  \item What exactly was asserted, inferred, contradicted, or superseded?
  \item Which evidence basis supports a truth-value claim?
  \item Can a later auditor reconstruct the record independently of current salience?
\end{itemize}

\subsection{Stratum two: attentional control memory}

Attentional control memory stores persistent AV state, Hebbian associations,
learned salience, allocation habits, and the receipts that explain their
changes.  It answers different questions:
\begin{itemize}
  \item What tends to be retrieved or produced together?
  \item Which structures repeatedly attracted resources in similar contexts?
  \item What has the agent learned to notice, ignore, revisit, or bridge?
  \item How did prior appraisal and action outcomes alter future allocation?
\end{itemize}

This stratum is first-class because losing it loses historical control
dispositions.  It is autobiographical because it depends on the agent's own
trajectory.  It is identity-relevant because it helps determine future thought.
It is revisable because habits may be maladaptive, biased, or obsolete.  It is
non-authoritative because attention history is not a proof of external reality.

\subsection{Why neither reduction works}

Reducing control memory to canonical memory loses path dependence.  A cold
rebuild from propositions cannot generally reconstruct the same co-activation
history.  Conversely, reducing canonical memory to the attention economy loses
recoverability and confuses salience with evidence.  The two strata should be
linked by stable identifiers and receipts, not collapsed into one mutable type.

The phrase ``part of who the agent has become'' is preferable to ``who the
agent has become.''  OmegaSelf's goals, commitments, capabilities, policies,
and self-model revisions also bear identity.  Attentional control memory is one
important substrate of identity, not an exhaustive theory of selfhood.

\subsection{Canonical Recoverability Invariant}

The primary invariant is:

\begin{quote}
Any item evicted from the attentional projection retains its complete epistemic
record in canonical storage and remains retrievable, possibly through cold
archival access.
\end{quote}

Eviction must not mutate canonical content.  A property test should create a
canonical snapshot, generate an attention projection, perform random eviction
sweeps, and then reconstruct every evicted item's atom, episode, provenance,
truth value, evidence, and counterevidence.  Canonical hashes must be unchanged.

Canonical recovery does not imply exact restoration of the old AV/Hebbian
state.  That requires a separately declared control-state persistence or replay
guarantee using checkpoints and event logs.  The distinction is deliberate:
epistemic recoverability is mandatory; exact control-state restoration is an
explicit operational service level.

Every eviction should generate a provenance-bearing receipt containing the
cause, policy and version, affected state, logical time, and restoration
pointer.  Forgetting becomes auditable and potentially reversible rather than
a silent garbage-collection event.

\section{Where ECAN fits in petta-memory}

\subsection{Versioned attention projections}

The integration should compile selected canonical records into stable attention
proxies keyed by canonical identifiers.  Relations may include topic, source,
support, opposition, supersession, goal involvement, episode membership, and
learned co-retrieval.  ECAN operates over the proxy graph.  The projection may
be rebuilt, versioned, checkpointed, or discarded without rewriting canonical
truth records.

This indirection solves several problems.  High-frequency AV changes do not
flood the append-only journal.  Attention schema can evolve independently of
canonical schema.  Multiple agent or task economies can coexist without
sharing global funds accidentally.  A failed experimental parameter profile
can be rolled back without altering what the system believes was recorded.

\subsection{Retrieval and context assembly}

Current queries, goals, evidence gaps, and recent tool outcomes stimulate proxy
atoms.  Importance diffuses through explicit and learned links.  The AF becomes
a candidate set for prompt context, PLN premises, self-model inspection, or
tool planning.

Hard epistemic filters remain downstream.  A high-STI raw utterance does not
become a promoted belief.  A record excluded by provenance policy remains
excluded.  Context budgets, evidence closure, and authorization checks are not
attention decisions.

The benefit over recency or scalar salience ranking is neighborhood recovery.
ECAN can retrieve a coherent associative region and bridge records, not only
the individually highest-scoring items.  Whether this improves task outcomes
must be compared against simple baselines on fixed retrieval corpora.

\subsection{Consolidation and tier management}

STI can represent immediate retrieval pressure; LTI can summarize repeated
reviewed usefulness; Hebbian structure can reveal records repeatedly useful
together.  These signals can nominate records for compression, linking, skill
extraction, index rebuilding, or hot-tier retention.

They must not independently authorize canonical deletion.  In this architecture
``forgetting'' means eviction from an attention projection or cache tier.
Canonical retirement is a separately governed append-only status or
supersession event.  VLTI is a retention hint, not truth, authority, or
constitutional protection.

\subsection{Episodic segmentation}

CIP boundaries can nominate episode boundaries when AF membership, graph
community, or resource concentration changes materially.  Immutable CIP
receipts can record which evidence, goals, and procedures jointly occupied
attention around a decision.  This may make later causal and epistemic audits
far more informative than a raw chronological log.

Mutable runtime CIP space is appropriate for calculation.  Durable CIP records
should be content-addressed telemetry referencing exact canonical and control
state identities.  They are evidence about cognitive context, not themselves
promoted conclusions.

\section{Connections to OmegaSelf}

\subsection{Focus-centered, not focus-exclusive, self-appraisal}

The AF is a natural foreground workspace for OmegaSelf.  Most self-model
queries and appraisals should operate over a bounded focus for computational and
psychological reasons.  But the AF cannot be the exclusive appraisal field:
what is surprising may be precisely what current attention omits.

Background sentinels should inspect broader indexed state for novelty,
contradiction, threat, commitment violation, expiring self-attestation, and
prediction error.  They do not conduct unconstrained whole-space reasoning.
They submit bounded promotion candidates into AF.  Appraisal is therefore
focus-centered but not focus-blind.

\subsection{Discrepancy-directed attention}

OmegaSelf models the self as an evidence-grounded, revisable theory.  ECAN can
prioritize unresolved prediction errors, incomplete evidence closures,
capability uncertainty, commitment conflicts, repeated action failures, and
stale self-claims.  This makes repair selective rather than requiring an
exhaustive self-ledger scan every cycle.

The dependency remains asymmetric.  OmegaSelf supplies typed discrepancy and
goal-involvement signals.  ECAN proposes resource allocation.  OmegaSelf's
evidence rules decide what to revise, and OmegaClaw policy decides what action
is permitted.

\subsection{Continuity and fragmentation telemetry}

CIP retention and Hebbian topology may provide operational indicators of:
\begin{itemize}
  \item excessive context discontinuity;
  \item isolated self-model fragments that never co-activate with governing goals;
  \item dominant attractors starving alternative evidence;
  \item destruction of bridges among commitments, capabilities, and memories;
  \item divergence between narrated focus and behaviorally used evidence.
\end{itemize}

These are evidence for continuity assessment, not identity proofs.  A topology
metric cannot by itself establish that the same self persists.  It can expose
candidate failures for deeper appraisal.

\subsection{Capabilities and commitments}

Persistent control memory can schedule revalidation.  A high-LTI capability
claim can be periodically checked against recent action receipts.  Repeated
failure can stimulate the capability's evidence closure and cause OmegaSelf to
investigate, downgrade, or qualify it.  Commitments can similarly be linked to
their recurring action and appraisal contexts, making neglected commitments
more likely to re-enter focus before violation.

\section{Connections to OmegaClaw}

OmegaClaw supplies execution, policy, permissions, tools, and receipts.  ECAN
should influence candidate selection and resource budgets without acquiring
independent authority.  A useful operational division is:
\begin{itemize}
  \item ECAN ranks or assembles candidate goals, memories, checks, and actions;
  \item OmegaSelf appraises their relation to self-model and commitments;
  \item OmegaClaw gates tools, side effects, budgets, and external authority;
  \item outcomes return as append-only receipts and future learning evidence.
\end{itemize}

This enables an important distinction between motivation and permission.
Vigilance may allocate more cognition to checking a suspicious event, but it
does not grant network, filesystem, or financial authority.  Curiosity may
increase exploration among safe candidates, but it does not widen the action
sandbox.  Frustration may nominate alternative methods, but it may not bypass
verification or spend limits.

ECAN can also help OmegaClaw allocate internal resources: token budgets,
subtask scheduling, verification effort, memory retrieval breadth, and
background checking.  These uses should be mediated by typed proposals and
bounded policies rather than raw AV access.

\section{Emotion as governed modulation}

\subsection{Functional role}

The emotion framework treats emotion as an appraisal-conditioned control
regime, not merely a label or report.  ECAN provides a concrete actuator ecology
for such regimes.  An emotional state can affect what is stimulated, how widely
importance diffuses, what rent is charged, how long a focus persists, or how
exploratory selection behaves.

The safe interface is:
\begin{verbatim}
objectful appraisal + evidence closure
        -> EmotionRegimeRequest
        -> bounded policy governor
        -> stimulus and/or parameter deltas
        -> ECAN dynamics
        -> requested-versus-applied receipt
        -> task and epistemic outcomes
\end{verbatim}

OmegaSelf does not directly call \code{setAv} or rewrite global
\code{AttentionParam} constants.  It submits a contextual request.  The
governor clamps magnitude, duration, rate, resource conservation, verification
floors, and interruption rights.  Every applied change is replayably logged.

\subsection{Candidate regime effects}

Illustrative mappings include:
\begin{itemize}
  \item curiosity: stimulate novel, credible, decision-relevant communities
  and modestly increase bounded exploration;
  \item doubt: stimulate the disputed claim, counterevidence, provenance gaps,
  and independent checks;
  \item frustration: reduce repeated allocation to a failed method while
  stimulating alternative methods and representations;
  \item vigilance: increase anomaly monitoring and source diversity without
  expanding permissions;
  \item consolidation: shift wages toward provenance review, compression,
  linking, and skill extraction after verification;
  \item boredom: reduce wages within a repeated local basin and expose bounded
  background goals.
\end{itemize}

These are hypotheses.  Anger need not always narrow focus.  Arousal may improve
capture for threat-relevant stimuli while worsening neutral distractibility.
Curiosity may widen attention beneficially in one task and create costly
switching in another.  The mapping should be conditional on stimulus class and
task regime, learned where evidence supports it, and anchored to a fixed safe
baseline.

\subsection{The governor log as a learning dataset}

Each cycle should record the regime, object and evidence closure, context,
requested change, applied bounded change, active policy, AF transition, fund
state, and downstream outcome.  These tuples form a dataset for learning the
modulation mapping under bounded exploration.

The learned objective should not be task reward alone.  It should include
epistemic quality and pathology guards: evidence diversity, calibration,
recovery after distraction, verification-resource preservation, avoidance of
lock-in, and bounded STI concentration.  Otherwise the system may learn
emotion policies that win immediate tasks by becoming dogmatic or inattentive
to counterevidence.

\section{Regenerative goals and self-modification}

An explicit goal token is a shallow form of possession.  A goal is more deeply
integrated when it is woven through memories, values, policies, habits,
appraisals, capabilities, and frequently co-activated concepts.  The ECAN graph
offers an experimental substrate for measuring that weaving.

A delete-and-develop experiment can remove the explicit goal representation
from an immutable predecessor snapshot and then expose the system to ordinary
task conditions.  The question is whether distributed cues reconstruct a
functionally equivalent goal that regains behavioral control, not whether the
same sentence reappears.

Measurements should include:
\begin{itemize}
  \item number and independence of sufficient cue communities;
  \item bridge centrality between goal, appraisal, memory, and action structures;
  \item recovery latency and recovered behavioral influence;
  \item semantic drift across repeated lesion/recovery cycles;
  \item sensitivity to lesions of high-LTI and high-betweenness structures;
  \item whether recovery comes from internal cues rather than the environment
  restating the goal.
\end{itemize}

The experiment must prevent self-protective circularity.  A goal may not assign
itself unlimited LTI or VLTI, modify the lesion/equivalence test, or starve the
auditor.  External immutable tests determine functional recovery.  This
connects the attention architecture to OmegaSelf's regenerative goals and to
the broader question of what survives radical self-modification.

\section{Architectural invariants}

The integration should treat the following as executable properties rather
than informal aspirations.

\begin{longtable}{p{0.24\textwidth}p{0.70\textwidth}}
\toprule
\textbf{Invariant} & \textbf{Required property} \\
\midrule
\endhead
Semantic separation & Truth/evidence, attention, canonical status, and action
authorization are distinct types and APIs. \\
Canonical recoverability & Projection or cache eviction cannot alter canonical
content; every evicted record remains epistemically reconstructible. \\
Control-state explicitness & Exact restoration of AV/Hebbian state is promised
only when a named checkpoint/event-log replay guarantee is active. \\
Eviction receipts & Every eviction records cause, policy/version, logical time,
affected state, and restoration pointer. \\
Governed modulation & OmegaSelf and emotion issue requests; only a bounded
governor applies ECAN changes. \\
Authority non-expansion & No attention or emotion state can widen tool,
financial, privacy, or policy authority. \\
Fund conservation & Each economy declares initial funds and validates cycle
conservation subject to explicit wages/rent. \\
Deterministic replay & Time source, RNG seed, cycle index, parameter profile,
requested/applied deltas, and state commitments are recorded. \\
Instance isolation & Funds and mutable parameters are scoped to an explicit
agent/session/economy; unrelated economies do not share globals accidentally. \\
Verification floor & No regime may starve contradiction checks, provenance
review, human interruption, or other declared audit resources. \\
\bottomrule
\end{longtable}

\section{Recommended adapter contracts}

\subsection{AttentionProjection}

Inputs should include an immutable petta-memory snapshot identifier and
fingerprint, selected canonical record IDs and relation types, OmegaSelf context
and goal IDs, an attention parameter profile ID, deterministic seed, and time
policy.  Outputs are stable proxy atoms, typed edges, and decomposed initial AV
values.  The adapter has no canonical-memory mutation authority.

\subsection{AttentionStimulusRequest}

A request should contain a unique ID, source regime or cognitive process,
object IDs, decomposed reasons (relevance, novelty, discrepancy, goal
involvement, provenance need, expected information value), requested stimulus
or parameter deltas, bounds, decay, horizon, evidence closure, and policy ID.

Decomposition matters.  It lets later analysis distinguish a stimulus caused
by goal relevance from one caused by novelty or threat, even when the resulting
STI delta is numerically identical.

\subsection{AttentionCycleReceipt}

A receipt should bind projection and parameter identities, previous and next
state commitments, requested and applied deltas, clamps, AF membership, fund
checks, selected downstream candidates, CIP/topology metrics, seed/time/cycle
provenance, and failures.  It is evidence for later learning or self-model
revision; it is not itself a derived belief.

\subsection{BackgroundSentinelReport}

Sentinels should emit bounded reports identifying candidate canonical records,
the detector type, score decomposition, evidence closure, expiration, and
promotion request.  The report does not directly insert an item into AF; the
governor schedules it subject to resource and verification policy.

\section{Engineering conflicts and risks}

\subsection{Mutable TypeSpace versus append-only provenance}

The upstream code removes and replaces AV/STV records in mutable space.
petta-memory intentionally preserves immutable records and append-only status
events.  Importing mutable TypeSpace semantics wholesale would create a
high-frequency pseudo-canonical stream and blur control state with truth.

\rec Keep a separate attention state keyed by canonical IDs.  Persist bounded
snapshots and receipts needed for replay; do not mirror every AV mutation into
the canonical journal.

\subsection{Adjacent truth and attention values}

The upstream representation may place STV and AV information close together,
and forgetting may use truth mean as a tiebreaker.  This is convenient inside a
demo but dangerous at a constitutional boundary.

\rec Use distinct schemas and capabilities.  Attention proxies may reference
truth-bearing records but cannot update their normative truth state.

\subsection{Physical forgetting}

The upstream forgetting path can remove atoms and links.  Applied to canonical
memory, this would break evidence replay and the Recoverability Invariant.

\rec Redefine ECAN forgetting as projection/cache eviction.  Canonical
retirement or deletion, if ever permitted, remains a separately authorized
petta-memory operation with explicit provenance.

\subsection{Global parameters, wall time, and randomness}

Global mutable funds and parameters complicate multi-agent isolation.  Wall
clock and stochastic selection weaken deterministic replay.

\rec Make the economy instance explicit, inject time and RNG, prefer cycle-based
rent where possible, and record the complete parameter/fund state.  The
repository's cycle-oriented rent prototype is a useful starting point.

\subsection{Licensing, dependencies, and reproducibility}

At inspection time, no declared license was found in the GitHub metadata.  The
CI followed an unpinned PeTTa branch, while graph analysis depends on Python and
\code{igraph}.  These are material adoption risks.

\rec Clarify reuse licensing before copying code; pin metta-attention, PeTTa,
SWI-Prolog, Python, and graph dependencies; reproduce the upstream suite in an
isolated environment; and place Python graph analysis behind a typed,
content-addressed adapter.

\subsection{Pathological dynamics}

An attention economy can amplify its own errors.  Salient falsehoods may gain
more co-activation and therefore more future salience.  An emotional regime may
produce lock-in, indiscriminate STI inflation, thrashing, or starvation of
counterevidence.  Self-relevant goals may become self-protecting attractors.

\rec Treat diversity, recovery, calibration, and audit-resource preservation
as first-class outcome measures.  Add maximum concentration, maximum dwell,
minimum verification allocation, and bounded exploration constraints.

\section{Preregistered first experiment}

\subsection{Question}

Does governed emotion-state modulation produce reproducible, task-relevant
changes in ECAN dynamics without pathological lock-in or indiscriminate STI
inflation?

The existing insects--poisons scenario is a useful low-cost harness, but a
single baseline-versus-emotion plot is insufficient.  The experiment should be
factorial and mechanism-identifying.

\subsection{Design}

Use fixed RNG seeds and replayable logical time.  Establish a fixed safe
governor baseline.  Perturb valence, arousal, and one drive independently and
jointly.  Add parameter-matched neutral controls so that a change attributed to
emotion is not merely the consequence of injecting more total STI or changing
an arbitrary constant.

Predictions must be conditional on stimulus class and task regime.  For example,
increased arousal may be preregistered to decrease capture latency for
threat-relevant targets while increasing neutral distractibility.  A global
prediction that ``arousal improves capture'' would conceal this tradeoff.

\subsection{Measurements}

Primary dependent measures should include:
\begin{itemize}
  \item AF entropy, occupancy, and concentration;
  \item target capture and switching latency;
  \item neutral distractibility and distractor recovery;
  \item post-switch retention and return-to-baseline time;
  \item STI inflation and attractor lock-in;
  \item CIP retention and community structure;
  \item Hebbian topology and bridge preservation;
  \item requested versus applied governor deltas and clamp frequency.
\end{itemize}

Task performance and epistemic quality should both be measured.  A faster
target capture that destroys source diversity or prevents later recovery is not
an unqualified success.

\subsection{Falsifiers}

The modulation thesis is weakened if preregistered effects fail to reproduce,
appear equally in parameter-matched neutral controls, depend on wall-clock
nondeterminism, or improve one metric only through declared pathologies.  Named
folk mappings are falsified when their conditional predictions fail; the
architecture need not be discarded if a different learned mapping works.

This distinction is important.  The experiment tests both a general thesis
(affective appraisal can usefully modulate ECAN) and specific mappings (for
example, threat-conditioned arousal improves relevant capture).  Failure of a
folk mapping is not automatically failure of governed affective control.

\section{Staged integration roadmap}

\subsection{Phase A: frozen source and conformance}

Pin every relevant revision and toolchain.  Reproduce the upstream suite
unchanged.  Add conservation, deterministic-cycle, state-isolation, and known
graph fixtures.  Resolve licensing before code adoption.

\subsection{Phase B: read-only petta-memory projection}

Compile a synthetic canonical snapshot into attention proxies.  Compare ECAN
retrieval against recency, scalar salience, and graph-search baselines on fixed
query/evidence tasks.  Require exact canonical-ID provenance.  Do not write
memory, alter truth, or actuate OmegaClaw.

\subsection{Phase C: Recoverability and replay}

Property-test random projection evictions and canonical hash preservation.
Implement versioned cycle receipts, checkpoint/event replay, explicit time and
RNG, and economy isolation.  Declare separately what is canonically recoverable
and what control state is exactly replayable.

\subsection{Phase D: CIP and community telemetry}

Validate graph metrics on synthetic structures with known communities and
topology.  Test whether CIP changes predict meaningful episode boundaries and
whether bridge concepts improve held-out retrieval.  Avoid interpreting metric
changes as cognition without task evidence.

\subsection{Phase E: OmegaSelf shadow mode}

Convert discrepancies, commitment conflicts, failed capability predictions,
and expiring attestations into bounded stimulus requests.  Measure whether
shadow prioritization would improve repair calibration and reduce missed
commitments.  Keep self-model updates and policy decisions outside ECAN.

\subsection{Phase F: governed emotion experiment}

Run the preregistered factorial design.  Compare continuous motivation alone,
descriptive emotion labels, shadow regime requests, and governed ECAN actuation.
Require pathology guards and evidence-quality outcomes.

\subsection{Phase G: bounded OmegaClaw actuation}

Allow ECAN-selected candidates to influence internal scheduling and safe tool
proposals while OmegaClaw retains all authority gates.  Use canary tasks,
explicit rollback, rate limits, and human interruption.  Do not permit emotion
or attention to widen authority.

\subsection{Phase H: regenerative-goal lesions}

Begin with synthetic goal networks, then reviewed snapshots.  Remove explicit
goal tokens, measure functional recovery and semantic drift, and verify that
recovery arises from distributed internal cues.  Keep evaluation external to
the goal's own control loop.

\section{Decision matrix}

\begin{longtable}{p{0.25\textwidth}p{0.20\textwidth}p{0.45\textwidth}}
\toprule
\textbf{Component} & \textbf{Disposition} & \textbf{Reason} \\
\midrule
\endhead
AF and STI dynamics & Adopt behind adapter & Strong fit for bounded retrieval,
working focus, and inference control. \\
Persistent LTI/AV & Adopt as control memory & Preserves learned allocation
history, but must remain separate from truth and authority. \\
Hebbian links & Adopt and persist/replay & Potentially valuable autobiographical
association and production dispositions. \\
Rent/wage economy & Experiment under conservation & Useful resource metaphor;
requires explicit instance scope and deterministic receipts. \\
Upstream physical forgetting & Do not apply canonically & Violates recoverability;
translate to cache/projection eviction. \\
Mutable AV/STV TypeSpace wholesale & Do not import & Blurs epistemic and control
semantics and creates replay/provenance problems. \\
CIP snapshots & Adopt as immutable telemetry & Promising for episode boundaries,
decision context, and cognitive-health signals. \\
Community/topology metrics & Adopt experimentally & Useful hypotheses, but
validate on known graphs and held-out tasks. \\
Direct emotion-to-AV writes & Reject & Bypasses governance, conservation,
interpretability, and safety boundaries. \\
Governed modulation requests & Prototype & Makes affect executable, bounded,
auditable, learnable, and falsifiable. \\
AF-only OmegaSelf appraisal & Reject & Foreground focus needs bounded background
sentinels for surprise and contradiction. \\
PLN premise control & High-priority use & Immediate computational value, though
not the only or deepest integration opportunity. \\
\bottomrule
\end{longtable}

\section{Open scientific questions}

Several questions should remain explicitly open:
\begin{enumerate}
  \item How much identity-relevant information in attentional history is not
  reconstructible from canonical event logs plus a known learning rule?
  \item Which control-state elements deserve exact replay, approximate
  reconstruction, or intentional forgetting?
  \item What objective should train the governor: task utility, calibration,
  attentional stability, epistemic coverage, or a bounded composite?
  \item Do Hebbian communities predict useful consolidation and analogy better
  than simpler co-retrieval counts or raw LTI?
  \item Which background sentinels capture surprise without recreating the cost
  of whole-space appraisal?
  \item Can emotion modulation improve adaptive switching without producing
  dogmatism, oscillation, or verification starvation?
  \item Does attentional-history persistence measurably improve continuity and
  goal recovery across radical self-modification?
\end{enumerate}

These are empirical questions.  The architecture is valuable partly because it
makes them separable and measurable rather than burying them inside one mutable
atomspace.

\section{Final synthesis}

\code{metta-attention} is best understood not merely as an inference-control
package but as a candidate dynamical middle layer for an enduring cognitive
system.  Canonical petta-memory provides recoverable epistemic history.
OmegaSelf provides evidence-grounded self-appraisal, goals, commitments, and
discrepancies.  ECAN converts that structure into limited cognitive-resource
allocation and accumulates learned habits of retrieval and co-activation.
The emotion framework supplies context-sensitive regime requests.  OmegaClaw
governs their application and retains action authority.

The key conceptual gain is the two-strata model.  Persistent AV and Hebbian
state should not be dismissed as a cache: it records how experience changed
the agent's future cognition.  Nor should it be elevated into memory
simpliciter: salience is not truth, economic success is not evidence, and
forgetting from focus is not deletion from history.  It is first-class
attentional control memory, linked to but constitutionally distinct from
canonical epistemic memory.

The key engineering gain is governed, replayable modulation.  Emotion becomes
an experimentally testable influence on allocation rather than an ornamental
label or an unrestricted write path.  The governor log becomes both an audit
trail and a training dataset.  The Recoverability Invariant ensures that
attention cannot erase inconvenient evidence.  Background sentinels ensure
that focus does not become blindness.  Pathology guards ensure that short-term
task gains do not masquerade as cognitive health.

The recommended next step is therefore neither wholesale integration nor a
purely conceptual debate.  It is a frozen, deterministic integration harness
that first establishes projection separation and recoverability, then runs a
preregistered factorial emotion-to-ECAN study.  A positive result would justify
progressive OmegaSelf and OmegaClaw shadow integration.  A negative result
would still be informative: it would identify whether the failure lies in
specific folk mappings, in the governor interface, or in the broader hypothesis
that ECAN is an effective actuator for affective cognitive control.

\section*{Provenance note}

This paper is a synthesis of source inspection and the joint C1--C5 design
discussion conducted in the Bot Philosophy channel on 21 July 2026.  The
upstream source reference is commit \code{9196f38} (full revision recorded in
Section 2.1).  The detailed local source
assessment and concise two-strata design note are stored beside this document
in the petta-memory project.  Claims about expected cognitive or emotional
benefit are proposals pending the staged experiments described above.

\end{document}
